% !TeX program = xelatex
% 编译：xelatex tensors_for_physics.tex （连续运行两次）
% 使用 TeX Live / MiKTeX 自带的 Fandol 中文字体；无需额外字体文件。
\documentclass[UTF8,11pt,a4paper,fontset=fandol]{ctexart}
\usepackage[left=22mm,right=22mm,top=21mm,bottom=22mm,headheight=14pt,headsep=7mm]{geometry}
\usepackage{amsmath,amssymb,mathtools,bm}
\usepackage{tikz-cd}
\usepackage{booktabs,tabularx,array}
\usepackage{enumitem}
\usepackage{needspace}
\AddToHook{cmd/subsubsection/before}{\Needspace{18\baselineskip}}
\usepackage{fancyhdr}
\usepackage[hidelinks,unicode,bookmarksnumbered=true]{hyperref}
\hypersetup{pdftitle={张量：从线性映射到指标计算},pdfsubject={面向物理系学生的张量入门},pdfauthor={}}
\ctexset{
 section={format=\large\bfseries,beforeskip=1.2ex plus .3ex minus .2ex,afterskip=.65ex},
 subsection={format=\normalsize\bfseries,beforeskip=.9ex plus .2ex minus .1ex,afterskip=.4ex},
 paragraph={format=\normalsize\bfseries,beforeskip=.6ex,afterskip=.5em}
}
\linespread{1.12}
\setlength{\parskip}{.15em}
\setlength{\parindent}{2em}
\setlength{\abovedisplayskip}{7pt plus 2pt minus 2pt}
\setlength{\belowdisplayskip}{7pt plus 2pt minus 2pt}
\setlength{\abovedisplayshortskip}{4pt plus 2pt minus 1pt}
\setlength{\belowdisplayshortskip}{5pt plus 2pt minus 1pt}
\setlist{nosep,leftmargin=2em,topsep=3pt}
\renewcommand{\arraystretch}{1.25}
\pagestyle{fancy}
\fancyhf{}
\fancyhead[L]{\small 张量：从线性映射到指标计算}
\fancyhead[R]{\small 物理数学基础}
\fancyfoot[C]{\small\thepage}
\renewcommand{\headrulewidth}{.3pt}
\fancypagestyle{plain}{\fancyhf{}\fancyfoot[C]{\small\thepage}\renewcommand{\headrulewidth}{0pt}}
\numberwithin{equation}{section}
\newcommand{\R}{\mathbb R}
\newcommand{\eps}{\varepsilon}
\newcommand{\dd}{\mathrm d}
\newcommand{\Tspace}[2]{\mathcal T^{#1}_{#2}(V)}
\DeclareMathOperator{\Hom}{Hom}
\DeclareMathOperator{\Bil}{Bil}
\DeclareMathOperator{\Mult}{Mult}
\DeclareMathOperator{\tr}{tr}
\DeclareMathOperator{\grad}{grad}
\newcommand{\key}[1]{\par\smallskip\noindent\textbf{#1}\par\smallskip}
\newcolumntype{Y}{>{\raggedright\arraybackslash}X}
\setlength{\emergencystretch}{2em}
\clubpenalty=10000
\widowpenalty=10000
\displaywidowpenalty=1000
\begin{document}
\thispagestyle{plain}
\begin{center}
  {\LARGE\bfseries 张量：从线性映射到指标计算\par}
\end{center}


\noindent\textbf{约定。}除另有说明外，$V$ 是 $d$ 维实线性空间，$m,n$ 是非负整数。$(m,n)$ 型张量由 $m$ 个 $V$ 因子和 $n$ 个 $V^*$ 因子构成，分量相应带 $m$ 个上指标、$n$ 个下指标。全文先讨论不依赖度规的结构，再引入度规。我们只考虑有限维的线性空间。



\section{线性空间}
\subsection{向量与线性空间}

所谓线性空间，是一个集合 $V$，其中的元素称为\textbf{向量}，并且在这些元素之间定义了两种运算：
\begin{itemize}
\item 向量加法：
\(    V\times V\longrightarrow V,
    (u,v)\longmapsto u+v;
   \)
\item 数乘：
\(    \mathbb{R}\times V\longrightarrow V,
    \qquad
    (\lambda,v)\longmapsto \lambda v.
   \)
\end{itemize}
这里暂时只考虑实线性空间，因此标量取自 $\mathbb{R}$。

这两种运算必须满足我们在线性代数中熟悉的一系列规则，例如
\begin{align}
u+v &= v+u,\\
(u+v)+w &= u+(v+w),\\
\lambda(u+v) &= \lambda u+\lambda v,\\
(\lambda+\mu)v &= \lambda v+\mu v,\\
(\lambda\mu)v &= \lambda(\mu v).
\end{align}
同时还要求存在零向量 $0\in V$ 和每个向量的加法逆元，并满足

$$
1v=v.
$$

这些条件的核心意义是：\textbf{向量可以做线性组合，而且线性组合仍然留在同一个空间中。}

因此，一个对象是不是向量，并不取决于它长得像不像一根箭头，而取决于它是否属于某个具有上述线性结构的集合。例如，

$$
\mathbb{R}^3
$$

中的三元数组当然可以组成线性空间；所有次数不超过 $n$ 的多项式也组成一个线性空间；某一类函数、矩阵，甚至微分方程的解，也都可能构成线性空间。

为了实际描述一个抽象的线性空间，我们通常选取一组\textbf{基}。

设

$$
\{e_1,\ldots,e_d\}
$$

是 $V$ 中的 $d$ 个向量。如果它们同时满足：

\begin{enumerate}
\item \textbf{线性无关}：若
\(    a^1e_1+\cdots+a^de_d=0,
   \)
则必有
\(    a^1=\cdots=a^d=0;
   \)

\item \textbf{张成整个空间}：任意 $v\in V$ 都可以写成
\[
v=v^1e_1+\cdots+v^de_d,
\]

\end{enumerate}
那么称 ${e_1,\ldots,e_d}$ 是 $V$ 的一组\textbf{基}。

“线性无关”保证这些基向量之间没有冗余，而“张成整个空间”保证它们又足够描述空间中的所有向量。因此，一旦选定一组基，任意向量 $v\in V$ 都能够\textbf{唯一}写成
\begin{equation}
v=\sum_{i=1}^{d}v^i e_i.
\label{eq:vexp}
\end{equation}

其中 $v$ 是向量，$e_i$ 是基向量，而 $v^i$ 是普通的实数，称为 $v$ 在这组基下的\textbf{分量}或\textbf{坐标}。

于是我们经常把这些分量排列成一个列向量

$$
[v]_{\{e_i\}}
=
\begin{pmatrix}
v^1\\
\vdots\\
v^d
\end{pmatrix}.
$$

这个列向量是 $v$ 在特定基 ${e_i}$ 下的一组坐标。

如果换成另一组基，$v$ 这个几何或代数对象并没有改变，但基向量 $e_i$ 会改变，相应的分量 $v^i$ 也会随之改变。二者的变化会恰好相互补偿，使得

$$
v=v^i e_i
$$

所表示的向量保持不变。

\section{1-form与对偶空间}
线性空间上的一个\textbf{线性函数}$T$是满足以下条件的，$V \rightarrow \mathbb{R}$（或者$\mathbb{C}$）的映射:
\begin{align}
    T(x+y) &= T(x)+T(y) \\
    T(\alpha x)&= \alpha T(x)，\alpha\in F
\end{align}
其中$F$是我们线性空间选择的数域，比如说$\mathbb{R}$或者$\mathbb{C}$，出于简便，我们接下来默认在$\mathbb{R}$上

可以证明，全部的线性函数构成了一个线性空间，比如说我们可以定义他们的加法和数乘：
\begin{align}
    (T_1+T_2)(x) &= T_1(x)+T_2(x) \\
    (\alpha T)(x) &= \alpha T(x)
\end{align}
我们把这个空间称作$V$的\textbf{对偶空间}，记为$V^*$。
对偶空间里的元素，称为\textbf{1-形式（1-form）}

自然出现了一个从 $V^*\times V\rightarrow\mathbb{R}$ 的双线性映射，称为\textbf{自然配对}：
\begin{equation}
    \langle T,x\rangle=T(x).
\end{equation}

这个配对还给出了一个从 $V$ 到双对偶空间 $V^{**}=(V^*)^*$ 的自然线性映射
\begin{equation}
    \iota:V\longrightarrow V^{**},\qquad
    \iota(x)(T)=T(x).
\end{equation}
也就是说，每个向量 $x\in V$ 都可以看成一个作用在 $V^*$ 上的线性函数：给它一个 $1$-形式 $T$，它返回实数 $T(x)$。

因为本文只考虑有限维空间，$\iota$ 实际上是一个同构：
\begin{equation}
    V\cong V^{**}.
\end{equation}
这个同构不依赖于基的选择，因此称为\textbf{自然同构}。读者大致可以认为$V^*$的对偶空间是$V$（在同构的意义下）
\section{张量积与多重线性函数}
\textbf{多重线性函数}指的是一个有若干自变量的函数，它对每个自变量分别都是线性的。例如，对第 $i$ 个自变量有
\begin{align}
     T(\ldots,x_i+y_i,\ldots,x_n)
     &=T(\ldots,x_i,\ldots,x_n)+T(\ldots,y_i,\ldots,x_n),\\
     T(\ldots,\alpha x_i,\ldots,x_n)
     &=\alpha T(\ldots,x_i,\ldots,x_n).
\end{align}
最常见的例子之一是行列式：把一个 $n\times n$ 矩阵看成由 $n$ 个列向量组成时，行列式就是关于这 $n$ 个列向量的 $n$ 重线性函数。

\textbf{张量积}用于将多重线性函数转为另一个线性空间里的线性函数。先从两个线性空间 $V,W$ 的情形开始。

我们希望构造一个新的线性空间 $V\otimes W$，并配备一个双线性映射
\begin{equation}
    \otimes:V\times W\longrightarrow V\otimes W,
    \qquad (v,w)\longmapsto v\otimes w,
\end{equation}
使得任意从 $V\times W$ 出发的双线性映射，都能够唯一地通过这个空间改写成线性映射。具体来说，对任意线性空间 $U$ 和任意双线性映射
\[
B:V\times W\longrightarrow U,
\]
都存在唯一的线性映射
\[
\widetilde B:V\otimes W\longrightarrow U,
\]
使得
\begin{equation}
    B=\widetilde B\circ\otimes,
    \qquad\text{即}\qquad
    B(v,w)=\widetilde B(v\otimes w).
\end{equation}
这个关系可以用下面的交换图表示：
\[
\begin{tikzcd}[column sep=large,row sep=large]
V\times W \arrow[r,"B"] \arrow[d,"\otimes"'] & U \\
V\otimes W \arrow[ur,dashed,"\exists!\,\widetilde B"'] &
\end{tikzcd}
\]
所谓\textbf{交换}，就是从左上角的 $(v,w)$ 到右上角的 $U$ 有两条路径：直接用 $B$，或者先送到 $v\otimes w$ 再用 $\widetilde B$；两条路径得到完全相同的结果。虚线上的 $\exists!$ 表示这样的线性映射 $\widetilde B$ 不仅存在，而且是唯一的。这一性质称为张量积的\textbf{普适性质}，它可以看作张量积与基无关的定义。

为了看见这个空间具体长什么样，设
\[
\{e_1,\ldots,e_d\}\quad\text{是 }V\text{ 的一组基},
\qquad
\{f_1,\ldots,f_r\}\quad\text{是 }W\text{ 的一组基}.
\]
那么 $V\otimes W$ 可以取为一个以
\[
e_i\otimes f_j,\qquad 1\le i\le d,\quad 1\le j\le r
\]
为基的 $dr$ 维线性空间。若
\[
v=v^i e_i,\qquad w=w^j f_j,
\]
就定义
\begin{equation}
    v\otimes w=v^i w^j\,(e_i\otimes f_j).
\end{equation}
从这个式子可以直接看出 $(v,w)\mapsto v\otimes w$ 对两个自变量分别是线性的，例如
\[
(v_1+v_2)\otimes w=v_1\otimes w+v_2\otimes w,
\qquad
v\otimes(w_1+w_2)=v\otimes w_1+v\otimes w_2.
\]
注意，$v\otimes w$ 这样的元素称为\textbf{简单张量}（或纯张量）；一般的 $V\otimes W$ 中的元素是若干简单张量的线性组合。

现在再看前面的双线性映射 $B$。因为它是双线性的，只要知道所有
\[
B(e_i,f_j)
\]
的值，就完全确定了 $B$。因此可以定义线性映射 $\widetilde B$ 满足
\begin{equation}
    \widetilde B(e_i\otimes f_j)=B(e_i,f_j),
\end{equation}
再由线性性延拓到整个 $V\otimes W$。于是
\[
\widetilde B(v\otimes w)
=v^i w^j\widetilde B(e_i\otimes f_j)
=v^i w^j B(e_i,f_j)
=B(v,w),
\]
这就具体展示了交换图为什么成立。由于 $e_i\otimes f_j$ 构成一组基，$\widetilde B$ 在这些基向量上的值已经被 $B$ 唯一确定，因此它也是唯一的。

特别地，当 $U=\mathbb{R}$ 时，任意双线性函数
\[
T:V\times W\longrightarrow\mathbb{R}
\]
都唯一对应于 $V\otimes W$ 上的一个线性函数
\[
\widetilde T:V\otimes W\longrightarrow\mathbb{R}.
\]
双线性函数等价于张量积空间上的普通线性函数。

完全相同的思想可以推广到多个线性空间。对
\[
T:V_1\times\cdots\times V_n\longrightarrow U
\]
这样的多重线性映射，都存在唯一线性映射
\[
\widetilde T:V_1\otimes\cdots\otimes V_n\longrightarrow U
\]
使得
\[
T(v_1,\ldots,v_n)
=\widetilde T(v_1\otimes\cdots\otimes v_n).
\]

\textbf{张量就是一类多重线性对象}，下一节我们会把这一点精确地写成 $(m,n)$ 型张量的定义。
\section{(m,n)型张量}
对一个线性空间$V$，我们可以在它上面定义(m,n)型张量，他是$V^* \times ...\times V^*\times V...\times V \rightarrow \mathbb{R}$上的多重线性函数，其中有$m$个$V^*$,$n$个$V$。

比如说，$V$中的一个矢量是(1,0)张量；$V^*$中的一个元素，也就是$V$上的线性函数，1-form，是(0,1)张量，标量，也就是$\mathbb{R}$上的一个数记作(0,0)张量。

我们也可以用 $V$ 和 $V^*$ 的张量积来生成高阶张量所处的线性空间。由于有限维情形下存在自然同构 $V\cong V^{**}$，所以 $(m,n)$ 型张量空间自然同构于
\[
V^{\otimes m}\otimes (V^*)^{\otimes n}.
\]
这里的“自然同构”强调：这种对应不需要额外选择一组基。
\section{张量代数}

对每个非负整数对 $(m,n)$，记
\[
\Tspace{m}{n}=V^{\otimes m}\otimes (V^*)^{\otimes n},
\]
并约定 $V^{\otimes 0}=(V^*)^{\otimes 0}=\R$。它与 $(m,n)$ 型多重线性函数空间自然同构，其中的元素就是 $(m,n)$ 型张量。

我们已经知道张量积可以用来生成更高阶的张量空间；在此基础上，我们可以预见应该存在一个方法，让两个张量可以生成更高阶的张量（因为张量积是有结合律的）。具体来说，两个张量可以作\textbf{张量积}：
\[
\otimes:\Tspace{m}{n}\times \Tspace{p}{q}\longrightarrow \Tspace{m+p}{n+q}.
\]
若 $S$ 是 $(m,n)$ 型张量，$T$ 是 $(p,q)$ 型张量，则 $S\otimes T$ 对前 $m$ 个 $V^*$ 输入和前 $n$ 个 $V$ 输入的作用来自 $S$，其余来自 $T$：
\[
\begin{aligned}
&(S\otimes T)(\omega^1,\ldots,\omega^m,\omega^{m+1},\ldots,\omega^{m+p},
v_1,\ldots,v_n,v_{n+1},\ldots,v_{n+q})\\
&=S(\omega^1,\ldots,\omega^m,v_1,\ldots,v_n)\,
T(\omega^{m+1},\ldots,\omega^{m+p},v_{n+1},\ldots,v_{n+q}).
\end{aligned}
\]
一般地，$S\otimes T\neq T\otimes S$；这不和我们前面提到的线性空间的张量积违背，我们前面已经给出了一个例子：线性空间的基经过张量积得到了张量空间里的基。

另一个基本运算是\textbf{缩并}。它把一个上指标和一个下指标配对，得到型数各减一的张量：
\[
C:\Tspace{m}{n}\longrightarrow \Tspace{m-1}{n-1}.
\]
抽象地说，缩并就是把一个 $V$ 因子和一个 $V^*$ 因子用自然配对
\[
\langle \omega,v\rangle=\omega(v)
\]
消去。缩并的定义不依赖基的选取。

把所有型数的张量空间直和起来，记
\[
\mathcal T(V)=\bigoplus_{m,n\ge 0}\Tspace{m}{n}.
\]
在张量积下，$\mathcal T(V)$ 构成一个分次代数，称为 $V$ 上的\textbf{张量代数}。它由 $V$ 与 $V^*$ 生成：$V$ 提供上指标，$V^*$ 提供下指标。

张量积提高型数，缩并降低型数；这是张量运算的两个最基本操作。

\section{分量，指标，爱因斯坦求和约定}

选定 $V$ 的一组基 $\{e_i\}_{i=1}^d$。对偶空间 $V^*$ 有唯一一组对偶基 $\{\eps^i\}$，满足
\[
\eps^i(e_j)=\delta^i_j.
\]
任意向量和 $1$-形式可写为
\[
v=v^i e_i,\qquad \omega=\omega_i\eps^i.
\]
这里 $v^i$ 是上指标分量，$\omega_i$ 是下指标分量。

$(m,n)$ 型张量 $T$ 可唯一展开为
\begin{equation}
T=T^{i_1\cdots i_m}_{j_1\cdots j_n}\,
e_{i_1}\otimes\cdots\otimes e_{i_m}\otimes
\eps^{j_1}\otimes\cdots\otimes \eps^{j_n}.
\end{equation}
其分量由对偶基作用给出：
\begin{equation}
T^{i_1\cdots i_m}_{j_1\cdots j_n}
=
T(\eps^{i_1},\ldots,\eps^{i_m},e_{j_1},\ldots,e_{j_n}).
\end{equation}
若 $\omega^a=\omega^a_{i_a}\eps^{i_a}$，$v_b=v_b^{j_b}e_{j_b}$，则
\begin{equation}
T(\omega^1,\ldots,\omega^m,v_1,\ldots,v_n)
=
T^{i_1\cdots i_m}_{j_1\cdots j_n}
\omega^1_{i_1}\cdots \omega^m_{i_m}
v_1^{j_1}\cdots v_n^{j_n}.
\end{equation}
这就是用分量计算张量的基本公式。

\textbf{爱因斯坦求和约定}：同一项中，一个上指标和一个下指标用相同字母表示时，自动对该字母从 $1$ 到 $d$ 求和：
\[
T^{i_1\cdots i_m}_{j_1\cdots j_n}\omega_{i_1}
:=
\sum_{i_1=1}^d T^{i_1\cdots i_m}_{j_1\cdots j_n}\omega_{i_1}.
\]
重复指标称为哑指标，可以成对改名；不重复的指标称为自由指标，等式两边必须一致。例如
\[
A^i{}_j B^j{}_k=C^i{}_k
\]
就是矩阵乘法，而
\[
A^i{}_i=\tr A.
\]
\textbf{在同一项中，不要让同一字母出现超过两次}。

张量积和缩并的分量形式分别为
\[
(S\otimes T)^{i_1\cdots i_m k_1\cdots k_p}_{j_1\cdots j_n l_1\cdots l_q}
=
S^{i_1\cdots i_m}_{j_1\cdots j_n}
T^{k_1\cdots k_p}_{l_1\cdots l_q},
\]
\[
(C T)^{i_1\cdots i_{m-1}}_{j_1\cdots j_{n-1}}
=
T^{i_1\cdots i_{r-1} k i_{r+1}\cdots i_m}_{j_1\cdots j_{s-1} k j_{s+1}\cdots j_n}.
\]
最后一式中 $k$ 是哑指标。

最后看基变换。设新基为
\[
e_{i'}=A^i{}_{i'}e_i.
\]
为保证同一个向量 $v=v^{i'}e_{i'}=v^i e_i$，分量必须满足
\[
v^{i'}=(A^{-1})^{i'}{}_i v^i,\qquad
\omega_{i'}=A^i{}_{i'}\omega_i.
\]
一般 $(m,n)$ 型张量分量中，每个上指标乘一个 $(A^{-1})$，每个下指标乘一个 $A$：
\[
T^{i'_1\cdots i'_m}_{j'_1\cdots j'_n}
=
(A^{-1})^{i'_1}{}_{i_1}\cdots (A^{-1})^{i'_m}{}_{i_m}
A^{j_1}{}_{j'_1}\cdots A^{j_n}{}_{j'_n}
T^{i_1\cdots i_m}_{j_1\cdots j_n}.
\]
因此，上指标与下指标的变换规律相反。正是这种“相反”，保证了张量作为整体对象在基变换下不变。注意：在没有度规时，上、下指标不能随意互换。

\section{度规诱导的同构，指标升降}

度规是 $V$ 上一个对称、非退化的双线性型
\[
g:V\times V\to\R,\qquad g(u,v)=g(v,u).
\]
按前面的约定，$g$ 是 $(0,2)$ 型张量，其分量为
\[
g_{ij}=g(e_i,e_j).
\]
非退化意味着矩阵 $(g_{ij})$ 可逆。记其逆矩阵为 $(g^{ij})$，满足
\[
g^{ik}g_{kj}=\delta^i_j.
\]
这里 $g^{ij}$ 是 $(2,0)$ 型张量的分量。

度规诱导出同构
\[
\flat:V\longrightarrow V^*,\qquad v^\flat(w)=g(v,w).
\]
非退化保证 $\flat$ 是单射；有限维且 $\dim V=\dim V^*$，所以它是同构。其逆记为
\[
\sharp:V^*\longrightarrow V.
\]
用分量写出就是
\[
(v^\flat)_i=g_{ij}v^j,\qquad
(\omega^\sharp)^i=g^{ij}\omega_j.
\]
这就是\textbf{指标下降}与\textbf{指标上升}。

对任意张量，可以用度规把上指标降下：
\[
T_{i_1}{}^{i_2\cdots i_m}{}_{j_1\cdots j_n}
=
g_{i_1 k}T^{k i_2\cdots i_m}{}_{j_1\cdots j_n},
\]
也可以把下指标升起：
\[
T^{i_1\cdots i_m j_1}{}_{j_2\cdots j_n}
=
g^{j_1 k}T^{i_1\cdots i_m}{}_{k j_2\cdots j_n}.
\]
多个指标同理，每次用一个 $g_{ij}$ 或 $g^{ij}$ 缩并即可。严格说，升降指标前后属于通过 $\flat,\sharp$ 同构起来的不同的张量空间；
物理中常说“把指标升上去”或“把指标降下来”，指的就是度规诱导的同构 $\flat,\sharp$ 以及相应的分量操作。
\section{例子}
\subsection{向量、余向量与缩并}

取 $V=\R^2$，基为 $\{e_1,e_2\}$，对偶基为
$\{\eps^1,\eps^2\}$。设
\[
v=2e_1-e_2,
\qquad
\omega=3\eps^1+4\eps^2.
\]
因此
\[
v^1=2,\quad v^2=-1,
\qquad
\omega_1=3,\quad \omega_2=4.
\]

向量与余向量的自然配对就是一次缩并：
\[
\begin{aligned}
\omega(v)
&=\omega_i v^i\\
&=\omega_1v^1+\omega_2v^2\\
&=3\times2+4\times(-1)\\
&=2.
\end{aligned}
\]
这里指标 $i$ 同时出现一次上标和一次下标，因此自动求和；
缩并以后不再留下指标，所得结果是一个标量。

再看一个 $(1,1)$ 型张量 $A$，设
\[
A^i{}_j=
\begin{pmatrix}
1&2\\
3&4
\end{pmatrix}.
\]
把它的上、下指标缩并，
\[
A^i{}_i
=A^1{}_1+A^2{}_2
=1+4
=5
=\tr A.
\]
因此，矩阵的迹正是 $(1,1)$ 型张量最基本的缩并。


\subsection{闵氏度规与指标升降}

取闵氏时空的号差为 $(-,+,+,+)$，并令 $c=1$。在惯性直角坐标中
\[
(g_{\mu\nu})
=
\begin{pmatrix}
-1&0&0&0\\
0&1&0&0\\
0&0&1&0\\
0&0&0&1
\end{pmatrix},
\qquad
(g^{\mu\nu})
=
\begin{pmatrix}
-1&0&0&0\\
0&1&0&0\\
0&0&1&0\\
0&0&0&1
\end{pmatrix}.
\]

设一个向量的逆变分量为
\[
V^\mu=(2,1,-3,4).
\]
下降指标意味着与度规缩并：
\[
V_\mu=g_{\mu\nu}V^\nu.
\]
逐个分量计算，
\[
\begin{aligned}
V_0&=g_{00}V^0=-1\times2=-2,\\
V_1&=g_{11}V^1=1,\\
V_2&=g_{22}V^2=-3,\\
V_3&=g_{33}V^3=4.
\end{aligned}
\]
因此
\[
V_\mu=(-2,1,-3,4).
\]
在这个坐标系和号差约定下，下降指标使时间分量变号，而空间分量不变。

反过来，若一个余向量为
\[
\omega_\mu=(5,-2,1,3),
\]
则升指标为
\[
\omega^\mu=g^{\mu\nu}\omega_\nu,
\]
所以
\[
\omega^\mu=(-5,-2,1,3).
\]

例如，向量与自身的闵氏内积为
\[
\begin{aligned}
V_\mu V^\mu
&=g_{\mu\nu}V^\mu V^\nu\\
&=-(V^0)^2+(V^1)^2+(V^2)^2+(V^3)^2\\
&=-4+1+9+16\\
&=22.
\end{aligned}
\]

\subsection{哑指标改名与化简}

哑指标只是求和变量，它使用什么字母并没有意义，因此可以成对改名。
例如
\[
A^i{}_jB^j{}_k-C^i{}_mB^m{}_k
\]
中，$i,k$ 是自由指标，而第一项中的 $j$ 和第二项中的 $m$
分别是哑指标。

把第二项中的哑指标 $m$ 改名为 $j$：
\[
C^i{}_mB^m{}_k
=
C^i{}_jB^j{}_k.
\]
于是原式可以写成
\[
\begin{aligned}
A^i{}_jB^j{}_k-C^i{}_mB^m{}_k
&=A^i{}_jB^j{}_k-C^i{}_jB^j{}_k\\
&=(A^i{}_j-C^i{}_j)B^j{}_k.
\end{aligned}
\]
所以在指标计算中，如果两个求和项看起来不能直接合并，
首先检查它们是否只是使用了不同的哑指标。

哑指标改名还常与张量的对称性一起使用。例如度规满足
\[
g_{ij}=g_{ji}.
\]
因此
\[
g_{ij}A^iB^j
=
g_{ji}A^iB^j.
\]
右边的 $i,j$ 都是哑指标，可以同时交换名称 $i\leftrightarrow j$：
\[
g_{ji}A^iB^j
=
g_{ij}A^jB^i
=
g_{ij}B^iA^j.
\]
从而
\[
g_{ij}A^iB^j
=
g_{ij}B^iA^j,
\]
这正是
\[
g(A,B)=g(B,A)
\]
在指标记号下的表现。

需要注意：只有哑指标可以任意改名；自由指标必须在等式两边保持一致。
改名时也不能让新的指标名称与同一项中已有的其他指标发生冲突。
\end{document}
